\documentclass[../../main.tex]{subfiles}
\begin{document}

\begin{center}
  {\fontsize{48pt}{58pt}\selectfont \textbf{1961}\normalsize\textbf{年度}}
\end{center}

\begin{table}[htb]
  \centering
  \caption{出題テーマと難易度}
  \begin{tabular}{|c|p{6.2cm}|c|c|c|}
    \hline
    \textbf{問題番号} & \textbf{テーマ} & \textbf{思考量} & \textbf{計算量} & \textbf{総合} \\
    \hline
    第1問 & 命題の否定・三角形の合同変換・四次関数の平行移動 & 普 & 普 & 普 \\
    \hline
    第2問 & 角の二等分線と外接円の接線 & 普 & 普 & 普 \\
    \hline
    第3問 & 三角形の角と正接の関係 & 普 & 普 & 普 \\
    \hline
    第4問 & 追跡運動と疲労の最小化 & 易 & 普 & 易 \\
    \hline
    第5問 & 絶対値の和の最小化 & 易 & 普 & 易 \\
    \hline
    第6問 & 速度が変化する2点の出会い & 普 & 普 & 普 \\
    \hline
  \end{tabular}
\end{table}

\begin{table}[htb]
  \centering
  \caption{解答}
  \renewcommand{\arraystretch}{1.4}
  \begin{tabular}{|c|c|p{9cm}|}
    \hline
    \textbf{問題番号} & \textbf{小問} & \textbf{答え} \\
    \hline
    第1問 & (i) & 「すべての実数に対し，そこで少なくとも一つの函数は0でない」 \\
    \cline{2-3}
    & (ii) & 平行移動・回転移動・必要なら線対称移動で重ねられる（証明問題） \\
    \cline{2-3}
    & (iii) & $h\ne0$では一致しない（証明問題） \\
    \hline
    第2問 & -- & $P$は$AD$の垂直二等分線上にある（証明問題） \\
    \hline
    第3問 & (1) & $\tan\alpha=\dfrac{2ar}{a^2-r^2}$ \\
    \cline{2-3}
    & (2) & $\tan\alpha\cos\theta=\tan\beta$（証明問題） \\
    \hline
    第4問 & -- & $B$の速さは毎時$2a\,\mathrm{km}$ \\
    \hline
    第5問 & -- & $n=50,51$のとき最小値$2500$ \\
    \hline
    第6問 & -- & $2$回 \\
    \hline
  \end{tabular}
\end{table}

\paragraph{試験形態}
6問．試験時間は要確認．

\end{document}
